\documentclass[11pt,a4paper]{article}

% ---------- Packages ----------
\usepackage[utf8]{inputenc}
\usepackage[T1]{fontenc}
\usepackage{lmodern}
\usepackage{amsmath,amssymb,siunitx,booktabs}
\usepackage{enumitem}
\usepackage{geometry}
\usepackage{fancyhdr}
\usepackage[most]{tcolorbox}   % enhanced, breakable
\usepackage{hyperref}

% ---------- Geometry ----------
\geometry{margin=2.5cm}

% ---------- Header/Footer ----------
\setlength{\headheight}{14pt}  % fancyhdr warning verhindern
\pagestyle{fancy}
\fancyhf{}
\lhead{Introduction to Electrical Engineering}
\rhead{Winter Term 2025}
\cfoot{\thepage}

% ---------- Exercise Environment ----------
\newcounter{exercise}
\newenvironment{exercise}[1][]
  {\refstepcounter{exercise}%
   \begin{tcolorbox}[enhanced, breakable, width=\textwidth,
     colback=blue!5!white,colframe=blue!60!black,
     fonttitle=\bfseries\large,
     title=Exercise~\theexercise:~#1]}
  {\end{tcolorbox}}

% ---------- Solution Environment ----------
\newif\ifshowsolutions
\showsolutionstrue  % comment to hide solutions

\newenvironment{solution}
  {\ifshowsolutions
     \begin{tcolorbox}[enhanced, breakable, width=\textwidth,
       colback=green!5!white,colframe=green!50!black,
       fonttitle=\bfseries\large,
       title=Solution]}
  {\end{tcolorbox}\fi}

% ---------- Document ----------
\begin{document}

% ----- Header Box -----
\begin{tcolorbox}[enhanced, breakable, colback=gray!10!white,
  colframe=gray!70!black, fonttitle=\bfseries\large,
  title={Exercise Sheet \#8}: AC Power and Resonances, sharp corners, width=\textwidth]
\begin{center}
  Discussion: Dec 15, 2025
\end{center}
\end{tcolorbox}
\vspace{1cm}

% ----- Exercise 1 -----
\begin{exercise}[...]
...
\end{exercise}

\ifshowsolutions
\begin{solution}
...
\end{solution}
\fi

% ----- Exercise 2 -----
\begin{exercise}[...]
...
\end{exercise}

\ifshowsolutions
\begin{solution}
...
\end{solution}
\fi

% ----- Exercise 3 -----
\begin{exercise}[...]
...
\end{exercise}

\ifshowsolutions
\begin{solution}
...
\end{solution}
\fi

% ----- Exercise 4 -----
\begin{exercise}[...]
...
\end{exercise}

\ifshowsolutions
\begin{solution}
...
\end{solution}
\fi

% ----- Exercise 5 -----
\begin{exercise}[...]
...
\end{exercise}

\ifshowsolutions
\begin{solution}
...
\end{solution}
\fi

% ----- Exercise 6 -----
\begin{exercise}[Series Resonance]
A series $RLC$ network has $R = 2\,\text{k}\Omega$, $L = 40\,\text{mH}$, and $C = 1\,\mu\text{F}$. Calculate the impedance at resonance and at one-fourth, one-half, twice, and four times the resonant frequency.
\end{exercise}

\ifshowsolutions
\begin{solution}
For s series RLC network with
\[
R=\SI{2}{\kilo\ohm},\qquad L=\SI{40}{\milli\henry},\qquad C=\SI{1}{\micro\farad}
\]
the series impedance is
\[
\underline{Z}(j\omega)=R+j\bigl(\omega L-\tfrac{1}{\omega C}\bigr).
\]

Resonance occurs when the net reactance is zero:
\[
\omega_0=\frac{1}{\sqrt{LC}}.
\]
Substitute the values:
\[
\omega_0=\frac{1}{\sqrt{(40\times10^{-3})(1\times10^{-6})}}=\SI{5000}{\radian\per\second}.
\]

For a given angular frequency \(\omega\),
\[
X_L=\omega L,\qquad X_C=\frac{1}{\omega C},\qquad X=X_L-X_C,
\]
\[
\underline{Z} = R + jX,\qquad |\underline{Z}| = \sqrt{R^2 + X^2},\qquad \varphi = \arctan\!\frac{X}{R}.
\]

We evaluate at \(\omega = k\omega_0\) for \(k=\tfrac14,\tfrac12,1,2,4\).

\bigskip
\begin{tabular}{@{}lcccccc@{}}
\toprule
$k$ & $\omega$ (rad/s) & $X_L$ (\si{\ohm}) & $X_C$ (\si{\ohm}) & $\underline{Z}$ (\si{k\ohm}) & $|\underline{Z}|$ (\si{k\ohm}) & $\varphi$ \\
\midrule
$0.25$ & $1250$ & $50$ & $800$ & $2 - j0.75$ & $2.136$ & $-20.6^\circ$ \\[4pt]
$0.5$ & $2500$ & $100$ & $400$ & $2 - j0.3$ & $2.022$ & $-8.5^\circ$ \\[4pt]
$1$ & $5000$ & $200$ & $200$ & $2$ & $2$ & $0^\circ$ \\[4pt]
$2$ & $10000$ & $400$ & $100$ & $2 + j0.3$ & $2.022$ & $+8.5^\circ$ \\[4pt]
$4$ & $20000$ & $800$ & $50$ & $2 + j0.75$ & $2.136$ & $+20.6^\circ$ \\
\bottomrule
\end{tabular}

\bigskip
\noindent Observations:
\begin{itemize}
  \item At resonance ($k=1$) the inductive and capacitive reactances cancel, leaving purely resistive impedance \(Z=R=\SI{2000}{\ohm}\).
  \item The network is capacitive ($X<0$) for frequencies below resonance and inductive ($X>0$) for frequencies above resonance.
  \item The magnitudes at pairs \(k=\tfrac14\) and \(k=4\) are equal (symmetry), and likewise for \(k=\tfrac12\) and \(k=2\).
\end{itemize}
\end{solution}
\fi

% ----- Exercise 7 -----
\begin{exercise}[Parallel Resonance]
Rework Exercise 6 if the elements are connected in parallel.
\end{exercise}

\ifshowsolutions
\begin{solution}
For the parallel connection the total admittance is
\[
\underline{Y}(j\omega) = \underbrace{\frac{1}{R}}_{G} + j\Bigl(\underbrace{\omega C}_{B_C} + \underbrace{\Bigl(-\frac{1}{\omega L}\Bigr)}_{B_L}\Bigr) = G + jB(\omega).
\]
The impedance is the reciprocal of the admittance:
\[
\underline{Z}(\omega) = \frac{1}{\underline{Y}(\omega)} = \frac{1}{G+jB} = \frac{G-jB}{G^2+B^2}.
\]
Its magnitude and phase are
\[
|\underline{Z}|=\frac{1}{\sqrt{G^2+B^2}},\qquad \varphi = \operatorname{atan2}\bigl(\Im\{Z\},\Re\{Z\}\bigr) = \arctan\!\frac{-B}{G}.
\]

Parallel resonance (net susceptance zero) occurs at the same angular frequency as series resonance:
\[
B(\omega_0) = 0 \quad\Rightarrow\quad \omega_0 = \frac{1}{\sqrt{LC}}.
\]
With the given values
\[
\omega_0 = \frac{1}{\sqrt{(40\times10^{-3})(1\times10^{-6})}} = \SI{5000}{\radian\per\second}.
\]
Also
\[
G = \frac{1}{R} = \frac{1}{2000} = \SI{5e-4}{\siemens}.
\]

At resonance $B(\omega_0)=0$, so
\[
Y(\omega_0) = G \qquad\Rightarrow\qquad Z(\omega_0) = \frac{1}{G} = R = \SI{2000}{\ohm}.
\]

We compute for each ratio \(k\):
\[
\omega = k\omega_0,\qquad B = \omega C-\frac{1}{\omega L},\qquad \underline{Z} = \frac{G-jB}{G^2+B^2}.
\]

\bigskip
\begin{tabular}{@{}lcccccc@{}}
\toprule
$k$ & $\omega$ (rad/s) & $B_C$ (\si{m\ohm}) & $B_L$ (\si{m\ohm}) & $\underline{Z}$ (\si{\ohm}) & $|\underline{Z}|$ (\si{\ohm}) & $\varphi$ \\
\midrule
$0.25$ & $1250$ & $1.25$ & $20$ & $1.42 + j\,53.30$ & $53$ & $+88.5^\circ$ \\[4pt]
$0.5$ & $2500$ & $2.5$ & $10$ & $8.85 + j\,132.74$ & $133$ & $+86.2^\circ$ \\[4pt]
$1$ & $5000$ & $5$ & $5$ & $2000$ & $2000$ & $0^\circ$ \\[4pt]
$2$ & $10000$ & $10$ & $2.5$ & $8.85 - j\,132.74$ & $133$ & $-86.2^\circ$ \\[4pt]
$4$ & $20000$ & $20$ & $1.25$ & $1.42 - j\,53.30$ & $53$ & $-88.5^\circ$ \\
\bottomrule
\end{tabular}

\bigskip
\noindent Notes on the results:
\begin{itemize}
  \item At resonance the susceptances of the inductor and capacitor cancel and the parallel branch behaves as a pure conductance $G$, giving $Z(\omega_0)=R$.
  \item For frequencies below resonance ($k<1$) the net susceptance $B$ is negative (inductive behavior), producing an impedance with a large positive imaginary part (inductive reactance).
  \item For frequencies above resonance ($k>1$) $B$ is positive (capacitive behavior), producing impedance with a large negative imaginary part (capacitive reactance).
  \item The impedances occur in conjugate pairs for reciprocal frequency ratios (e.g. $k=0.25$ and $k=4$; $k=0.5$ and $k=2$).
\end{itemize}
\end{solution}
\fi

\end{document}